不完备性定理的证明要求我们找到一种方法，能够用该理论自身的语言（即算术的语言）来谈论某个理论（例如 $\Th{PA}$）内部的可证性。但算术语言只处理数字，不处理公式或推导。解决这一问题的方案是定义一个从公式和推导到数字的系统映射。与公式或推导相对应的数字称为它的\textbf{Gödel数}。若 $!A$ 是一个公式，则 $\Gn{!A}$ 即为其Gödel数。我们证明了，对公式的重要操作会转化为关于相应Gödel数的原始递归函数。例如，用项~$t$ 替代~$!A$ 中变元~$x$ 的所有自由出现这一操作 $\Subst{!A}{t}{x}$，对应着一个算术函数 $\fn{subst}(n,m,k)$；将该函数应用于 $!A$、$t$ 和 $x$ 的Gödel数时，便得到~$\Subst{!A}{t}{x}$ 的Gödel数。换言之，$\fn{subst}(\Gn{!A}, \Gn{t}, \Gn{x}) = \Gn{\Subst{!A}{t}{x}}$。
类似地，推导的性质会转化为关于相应Gödel数的原始递归关系。特别地，性质 $\fn{Deriv}(n)$（当 $n$ 是自然演绎中一个正确推导的Gödel数时成立）是原始递归的。证明这些性质是原始递归的，需要相当多的工作，有时还需要一些技巧，并且本质上依赖于“处理序列是原始递归的”这一事实。如果理论~$\Th{T}$ 是可判定的，那么我们可以用 $\fn{Deriv}$ 来定义一个可判定的关系 $\Prf[\Th{T}](n, m)$，当 $n$ 是从~$\Th{T}$ 推出Gödel数为~$m$ 的语句的推导的Gödel数时，该关系成立。如果~$\Th{T}$ 的公理集合是原始递归的，那么这个关系就是原始递归的；而如果~$\Th{T}$ 的公理只是可判定的但非原始递归的，那么这个关系就仅是一般递归的。